\documentclass[12pt,a4paper]{article}
\usepackage[spanish,es-tabla]{babel}
\usepackage[utf8]{inputenc}
\usepackage[T1]{fontenc}
\usepackage{graphicx}
\usepackage{listings}
\usepackage{xcolor}
\usepackage{geometry}
\usepackage{amsmath}
\usepackage{array}
\usepackage{booktabs}
\usepackage{float}
\usepackage{longtable}
\usepackage{fancyhdr}
\usepackage{hyperref}
\usepackage{enumitem}

\geometry{margin=2.5cm}

\definecolor{codegray}{rgb}{0.97,0.97,0.97}
\definecolor{commentcolor}{rgb}{0,0.5,0}
\definecolor{keywordcolor}{rgb}{0,0,0.6}
\definecolor{stringcolor}{rgb}{0.6,0,0}
\definecolor{numbercolor}{rgb}{0.5,0.5,0.5}

\lstdefinestyle{codestyle}{
    backgroundcolor=\color{codegray},
    commentstyle=\color{commentcolor}\itshape,
    keywordstyle=\color{keywordcolor}\bfseries,
    numberstyle=\tiny\color{numbercolor},
    stringstyle=\color{stringcolor},
    basicstyle=\ttfamily\small,
    breakatwhitespace=false,
    breaklines=true,
    captionpos=b,
    keepspaces=true,
    numbers=left,
    numbersep=5pt,
    showspaces=false,
    showstringspaces=false,
    showtabs=false,
    tabsize=2,
    language=Python
}

\lstdefinestyle{grammartyle}{
    backgroundcolor=\color{codegray},
    basicstyle=\ttfamily\small,
    breaklines=true,
    numbers=none,
    showstringspaces=false,
    tabsize=2
}

\lstdefinestyle{bashstyle}{
    backgroundcolor=\color{codegray},
    basicstyle=\ttfamily\small,
    breaklines=true,
    numbers=none,
    showstringspaces=false,
    tabsize=2,
    language=bash
}

\pagestyle{fancy}
\fancyhf{}
\fancyhead[L]{Proyecto: Analizador Sintáctico LL(1) - SmartRule}
\fancyhead[R]{Teoría de la Computación - INFO1148}
\fancyfoot[C]{\thepage}

\begin{document}

% ===================== PORTADA =====================
\begin{titlepage}
    \centering
    \vspace*{-1cm}
    {\Large \textbf{Universidad Católica de Temuco}\par}
    {\large \textbf{Facultad de Ingeniería, Escuela de Ingeniería Informática}\par}
    \vspace{1.5cm}
    
    {\LARGE \textbf{INFORME DEL PROYECTO}\par}
    \vspace{0.5cm}
    {\Large \textbf{Diseño e Implementación de un Analizador Sintáctico LL(1)}\par}
    \vspace{1cm}
    
    {\Large \textbf{Tema:} SmartRule - DSL para Automatización IoT y Domótica\par}
    \vspace{2cm}
    
    \begin{center}
    \begin{tabular}{|p{4.5cm}|p{9cm}|}
    \hline
    \textbf{Dato} & \textbf{Información} \\ 
    \hline
    Código curso & INFO1148 \\ 
    \hline
    Integrantes & José Sáez (Jefe de grupo) \newline Catalina Vergara \newline Alexa Galeano \\ 
    \hline
    Profesor & Prof. M. Lévano \\ 
    \hline
    Fecha de entrega & % TODO: confirmar fecha de entrega en Blackboard \\ 
    \hline
    \end{tabular}
    \end{center}
    
    \vfill
    
    {\large Temuco, Chile\par}
    {\large \today\par}
\end{titlepage}

% ===================== INDICE =====================
\newpage
\tableofcontents
\newpage
\listoftables
\newpage
\listoffigures
\newpage

% ===================== 1. INTRODUCCIÓN =====================
\section{Introducción}

\subsection{Antecedentes}
La teoría de autómatas y los lenguajes formales constituyen la base fundamental para el diseño y construcción de compiladores e intérpretes. Dentro de este proceso, el análisis sintáctico (también conocido como \textit{parsing}) tiene como finalidad verificar que una secuencia de tokens generada por el analizador léxico, cumpla con la estructura gramatical definida para un lenguaje determinado. 

Para lograrlo, se utilizan Gramáticas Libres de Contexto (GLC), las cuales permiten describir de manera formal la sintaxis de lenguajes de programación, así como Lenguajes de Dominio Específico (\textit{Domain Specific Languages - DSL}).

\subsection{Planteamiento del Problema}
Se requiere diseñar e implementar un analizador sintáctico capaz de reconocer, validar y representar sintácticamente programas escritos en un Lenguaje de Dominio Específico creado por el grupo. Además, debe ser capaz de detectar errores sintácticos, entregar mensajes claros y construir una representación estructurada de la entrada mediante un Árbol Sintáctico Abstracto (AST).

\subsection{Objetivo General}
Diseñar e implementar un analizador sintáctico LL(1) descendente recursivo para el lenguaje de dominio específico \textit{SmartRule}, aplicando Gramáticas Libres de Contexto, conjuntos \textit{FIRST} y \textit{FOLLOW}, análisis LL(1), construcción del Árbol Sintáctico Abstracto (AST) y detección de errores sintácticos.

\subsection{Objetivos Específicos}
\begin{itemize}[leftmargin=1cm]
    \item Definir el dominio, propósito, tokens, palabras reservadas y sintaxis del lenguaje \textit{SmartRule}.
    \item Diseñar una Gramática Libre de Contexto (GLC) correcta, no ambigua, completa y coherente con el lenguaje propuesto.
    \item Analizar la gramática, identificando y corrigiendo recursividad, factorización y posibles ambigüedades.
    \item Calcular y justificar los conjuntos \textit{FIRST} y \textit{FOLLOW} para cada no-terminal.
    \item Construir la tabla predictiva LL(1), verificando la ausencia de conflictos.
    \item Implementar un analizador léxico y un analizador sintáctico LL(1) descendente recursivo, desarrollados íntegramente por el grupo.
    \item Construir un Árbol Sintáctico Abstracto (AST) que represente correctamente la estructura del programa.
    \item Implementar un mecanismo de detección y reporte de errores léxicos y sintácticos con mensajes claros y localizados.
    \item Desarrollar un conjunto de pruebas exhaustivo, incluyendo casos válidos, inválidos, casos límite y una extensión creativa integrada a la gramática.
    \item Documentar todo el proceso mediante un informe técnico en \LaTeX\ y un \texttt{README.md} reproducible.
\end{itemize}

\newpage
